\slajdid{06-s023}
\begin{frame}[label=06-s023]{Zaštitni kodovi – CRC – Cyclic Redundancy Check}
\gusttekst Predstava binarne vrednosti u obliku polinoma
\[b_{n-1}b_{n-2}b_{n-3}\ldots b_1b_0\equiv b(X)=b_{n-1}X^{n-1}+b_{n-2}X^{n-2}+b_{n-3}X^{n-3}+\ldots+b_1X^1+b_0X^0\]
\begin{center}\begin{tabular}{c*{10}{c}}\toprule
$\mathbf{B}$ & 1 & 0 & 1 & 0 & 1 & 0 & 0 & 1 & 0 & 1\\\midrule
$\mathbf{P}$ & $X^9$ & 0 & $X^7$ & 0 & $X^5$ & 0 & 0 & $X^2$ & 0 & $X^0{=}1$\\\bottomrule\end{tabular}
\end{center}
\[1010100101\equiv X^9+X^7+X^5+X^2+1\]
\begin{columns}[c,onlytextwidth]\column{.30\textwidth}Osnova CRC kodova je\\Generator polynomial stepena $k$
\column{.68\textwidth}\[g(X)=g_kX^k+g_{k-1}X^{k-1}+g_{k-2}X^{k-2}+\ldots+g_1X^1+g_0X^0\]\end{columns}
\vspace{1.5mm}\centerline{Na osnovu generator polinoma dodaje se $k$ zaštitnih bita}\vspace{1.5mm}
\begin{center}Vrednost zaštitnih bita se određuje na osnovu ostatka koji se dobija deljenjem
\[\left(P\times X^k\right)\div g(X)\]
{\Large\color{DEcrvena}po modulu 2\par}\end{center}
\end{frame}
